\documentclass[10pt,a4paper]{amsart}
\usepackage[margin=19mm]{geometry}
\usepackage{amsmath,amssymb,mathtools}
\usepackage[scr=rsfs,cal=euler]{mathalfa}
\usepackage{xfrac}
\usepackage[unicode,hidelinks,bookmarks=false]{hyperref}
\newcommand{\ie}{i.e.\ }
\newcommand{\cf}{cf.\ }
\newcommand{\resp}{respectively\ }
\newcommand{\Subsection}{\subsection}
\providecommand{\nobreakdash}{\nobreak\hbox{-}}
\setlength{\emergencystretch}{2em}
\multlinegap0pt
\usepackage{amssymb}
\usepackage{mathtools}
\usepackage{xcolor}
\newtheorem{conjecture}{Conjecture}
\newtheorem{theorem}{Theorem}
\newtheorem{lemma}{Lemma}
\newtheorem{definition}{Definition}
\newcommand{\K}{\mathbb K}
\newcommand{\M}{\mathfrak M}
\newcommand{\Z}{\mathbb Z}
\newcommand{\cV}{\mathcal V}
\newcommand{\eps}{\varepsilon}
\title{Topological string blowup equations via stable pairs}
\author{Lutian Zhao}
\date{Source: arXiv:2608.25397v2 (CC BY 4.0)}
\begin{document}
\maketitle

\noindent\emph{Source and licence.} Topological string blowup equations via stable pairs. Version arXiv:2608.25397v2 (CC BY 4.0), distributed by arXiv under the Creative Commons Attribution 4.0 International licence, http://creativecommons.org/licenses/by/4.0/, as recorded in the arXiv licence metadata for this exact version and shown on the abstract page https://arxiv.org/abs/2608.25397v2. Full licence notice: \texttt{licenses/arxiv-2608.25397-CC-BY-4.0.txt}.\par\smallskip

\noindent\textit{Editorial scope.} Counted unit: the article's lemma lem:bilateral-u-zero-limit (source0.tex lines 4160--4167). The article's complete original proof of it is delivered with it (source0.tex lines 4169--4205, one \texttt{proof} environment of 37 lines). No mathematics is written, repaired or re-derived: the statement and the proof are the article's own lines, in the article's own notation, and no retained line is substituted or re-flowed.  Retained uncounted as the source-local context the proof needs: lines 2180--2184; lines 2234--2242; lines 2568--2600; lines 2603--2643; lines 2673--2703; lines 2718--2724; lines 2765--2785; lines 3160--3166; lines 3911--3915; lines 3917--3931; lines 3940--3962; lines 4011--4016; lines 4082--4101; lines 4146--4149.  Nothing was omitted from any retained span, so the package carries no omission record.  No retained line contains a citation, so the wrapper carries no bibliography and no bibliography entry.  No retained line contains a figure environment, an image-inclusion command or drawing code, so no image is delivered and the package declares no asset dependency.  Editorial formatting only, in addition: an AMS \texttt{amsart} preamble that restates the article's own declarations and macros used by the retained lines, namely the environments \texttt{conjecture}, \texttt{theorem}, \texttt{lemma}, \texttt{definition} on separate counters, together with the standard packages those lines need.  The retained environments print in this document as Conjecture 1, Theorem 1, Lemma 1, Definition 1 in order of appearance, whereas the article prints the same items with the numbering of its own full document.  The complete native source of the article as distributed in the arXiv e-print for this exact version is bundled unchanged as \texttt{sources/208577/source0.tex}.
\begin{equation}
 (w_1,w_2,w_3)
 =(\xi_1ab,\xi_2a/b,\xi_3a^{-2}).
\label{eq:P2-iterated-cocharacter}
\end{equation}

\begin{conjecture}
\label{conj:P2-vertex-comparison}
For the iterated limit of the subgroup \eqref{eq:P2-iterated-cocharacter}, the coefficientwise slope limit of the symmetrized stable pair localization formula is
\begin{equation}
 Z_{X}^{\mathrm{PT},\mathrm{ref}}(Q;t,q)
 =\mathsf Z_{\mathbb P^2}^{\mathrm{vert}}(Q;t,q).
\label{eq:P2-PT-equals-IK-vertex}
\end{equation}
\end{conjecture}

Write the framing weights symmetrically,
\[
 q_i=e^{-\eps_i},\qquad u=v^2,\qquad
 (e^{a_1},e^{a_2})=(v,v^{-1}),
\]
and define \(\mathfrak p\) by
\begin{equation}
 \Lambda^4(q_1q_2)^{3/2}=v\mathfrak p.
\label{eq:NY-p-variable}
\end{equation}
Multiplying both framing weights by \(v\) changes them from \((v,v^{-1})\) to \((u,1)\); it leaves the tangent character unchanged and multiplies \(\det\cV\) on \(\M(2,N)\) by \(v^N\).  Let
\begin{equation}
\begin{aligned}
 \Lambda_{\mathrm c}(\mathfrak p,u;q_1,q_2)
 &:={}
 \left(\frac{u^{1/2}\mathfrak p}{(q_1q_2)^{3/2}}\right)^{1/4},\\
 \mathscr Z(\mathfrak p,u;q_1,q_2)
 &:={}
 Z_1\left(
 -\log q_1,-\log q_2,
 \left(\tfrac12\log u,-\tfrac12\log u\right);
 \Lambda_{\mathrm c}(\mathfrak p,u;q_1,q_2)
 \right).
\end{aligned}
\label{eq:NY-mathscrZ-definition}
\end{equation}
We regard \(\eps_i\) and \(\log u\) as formal additive parameters, or fix analytic branches, and choose the fourth root in \(\Lambda_{\mathrm c}\) once and for all.  By the change of framing, the framed sheaf series of \(\mathscr Z(\mathfrak p,u;q_1,q_2)\) is
\begin{equation}
 \sum_{N\geq0}\mathfrak p^N
 \chi_{\mathbf T_{\mathrm{fr}}}^{\mathrm{loc}}\bigl(\M(2,N),\det\cV\bigr)
\label{eq:NY-centered-ADHM}
\end{equation}
with framing weights \((u,1)\).

\begin{theorem}
\label{thm:NY-native}
For \(k\in\{0,1\}\) and \(d_{\mathrm{NY}}\in\Z_{\geq0}\), let
\[
 c_{k,d_{\mathrm{NY}}}
 =\frac{k^3-k-2d_{\mathrm{NY}}+2}{24},\qquad
 \ell_n=(n-k/2,-n+k/2),
\]
and
\begin{equation}
\begin{aligned}
 \widehat Z_{1,k,d_{\mathrm{NY}}}
 (\eps_1,\eps_2,\vec a;\Lambda)
 ={}&e^{c_{k,d_{\mathrm{NY}}}s_\eps}
 \sum_{n\in\Z}
 Z_1\!\left(
 \begin{gathered}
 \eps_1,\eps_2-\eps_1,\vec a+\eps_1\ell_n;\\[-2pt]
 \Lambda e^{\eps_1(2d_{\mathrm{NY}}+k-3)/8}
 \end{gathered}\right)\\
 &\hspace{22mm}\times
 Z_1\!\left(
 \begin{gathered}
 \eps_1-\eps_2,\eps_2,\vec a+\eps_2\ell_n;\\[-2pt]
 \Lambda e^{\eps_2(2d_{\mathrm{NY}}+k-3)/8}
 \end{gathered}\right).
\end{aligned}
\label{eq:NY-native-additive}
\end{equation}
Then
\begin{equation}
\begin{aligned}
 \widehat Z_{1,0,d_{\mathrm{NY}}}
 (\eps_1,\eps_2,\vec a;\Lambda)
 &=Z_1(\eps_1,\eps_2,\vec a;\Lambda),
 &&0\leq d_{\mathrm{NY}}\leq2,\\
 \widehat Z_{1,1,1}(\eps_1,\eps_2,\vec a;\Lambda)&=0.
\end{aligned}
\label{eq:NY-unity-vanishing}
\end{equation}
\end{theorem}

In multiplicative form: for
\[
 (k,d_{\mathrm{NY}})\in
 \{(0,0),(0,1),(0,2),(1,1)\},
\]
and in the notation of \eqref{eq:NY-mathscrZ-definition},
\begin{equation}
\begin{aligned}
 &\sum_{n\in\Z}
 \frac{
 \mathscr Z(\mathfrak p q_1^{n-d_{\mathrm{NY}}-k},u q_1^{k-2n};
 q_1,q_2/q_1)}
 {\mathscr Z(\mathfrak p,u;q_1,q_2)}\\[-1mm]
 &\hspace{18mm}\times
 \mathscr Z(\mathfrak p q_2^{n-d_{\mathrm{NY}}-k},u q_2^{k-2n};
 q_1/q_2,q_2)\\
 &\qquad=
 \begin{cases}
 (q_1q_2)^{(1-d_{\mathrm{NY}})/12},
   &k=0,\ d_{\mathrm{NY}}=0,1,2,\\
 0,&(k,d_{\mathrm{NY}})=(1,1).
 \end{cases}
\end{aligned}
\label{eq:NY-multiplicative}
\end{equation}
Indeed, by \eqref{eq:NY-p-variable} on each chart, the shifts of \(\Lambda\) and \(\vec a\) in \eqref{eq:NY-native-additive} become
\[
 \mathfrak p_i=\mathfrak p q_i^{n-d_{\mathrm{NY}}-k},\qquad
 u_i=uq_i^{k-2n}.
\]
Divide \eqref{eq:NY-native-additive} by \(Z_1\) and use \eqref{eq:NY-unity-vanishing}; for \(k=0\), \(e^{-c_{0,d_{\mathrm{NY}}}s_\eps} =(q_1q_2)^{(1-d_{\mathrm{NY}})/12}\).

\begin{equation}
 F_{\mathbb P^2}^{(0)}(T;\alpha,\beta)
 =\frac{T^3}{18\alpha\beta}
 -\frac{\pi^2T}{6\alpha\beta}+\frac{T}{12}
 +\frac{(\alpha+\beta)^2T}{24\alpha\beta}.
\label{eq:P2-degree-zero-polynomial}
\end{equation}

\begin{lemma}
\label{lem:NY-chart-substitutions}
Take \((k,d_{\mathrm{NY}})\in\{(0,1),(0,2),(1,1)\}\) and replace \(n\) by \(-n\) in \eqref{eq:NY-multiplicative}.  On chart \(i=1,2\),
\begin{equation}
 \mathfrak p_i=\mathfrak p q_i^{-n-d_{\mathrm{NY}}-k},\qquad
 u_i=uq_i^{2n+k},\qquad
 \mathfrak a_i=\frac{\mathfrak a}{q_i},\qquad
 z_i=zq_i^{-n-d_{\mathrm{NY}}-k+3/2}.
\label{eq:NY-chart-substitutions}
\end{equation}
With \(Q_i=-z_i u_i^2\),
\begin{equation}
 \frac{Q_i}{Q}
 =q_i^{3n-d_{\mathrm{NY}}+k+3/2}
 =q_i^{R_{r,n}},
 \qquad
 R_{r,n}=3n+\frac r2,\qquad r=3-2d_{\mathrm{NY}}+2k.
\label{eq:NY-plane-shift}
\end{equation}
The three cases give \(r=1,-1,3\) respectively.
\end{lemma}

\begin{conjecture}
\label{conj:P2-asymptotic-resummation}
The analytic hypotheses \textup{(A1)--(A5)} hold for each of
\[
 (k,d_{\mathrm{NY}})=(0,1),\ (0,2),\ (1,1).
\]
\end{conjecture}

\begin{equation}
 v(\gamma)\geq-Ch(\gamma)
 \qquad(\gamma\in\Gamma_{\mathfrak S}).
\label{eq:semigroup-cone-bound}
\end{equation}

\begin{definition}
\label{def:completed-semigroup-algebra}
Let \(\mathscr H_{\mathfrak S}\) be the set of formal sums
\begin{equation}
 \sum_{\gamma\in\Gamma_{\mathfrak S}}c_\gamma \mathbf e^\gamma,
 \qquad c_\gamma\in\K_{\eps^\circ},
\label{eq:completed-semigroup-algebra}
\end{equation}
such that for every \(H,V\in\mathbb R\) only finitely many nonzero terms satisfy \(h(\gamma)\leq H\) and \(v(\gamma)\leq V\).  For \(H\geq0\), let
\[
 F^{>H}\mathscr H_{\mathfrak S}
 =\left\{\sum_\gamma c_\gamma \mathbf e^\gamma:
 c_\gamma=0\text{ whenever }h(\gamma)\leq H\right\}.
\]
\end{definition}

\begin{lemma}
\label{lem:completed-semigroup-support}
Multiplication is well defined in \(\mathscr H_{\mathfrak S}\), each \(F^{>H}\mathscr H_{\mathfrak S}\) is an ideal, and the algebra is complete and separated.  Every element of \(1+\mathfrak m_A\) is invertible, and \(\operatorname{CT}_A\) is a continuous ring homomorphism.  For \(i\in\{1,2\}\), \(n\in\Z\), and
\[
 (k,d_{\mathrm{NY}})\in\{(0,1),(0,2),(1,1)\},
 \qquad r=3-2d_{\mathrm{NY}}+2k,
\]
the chart translations
\[
 A\longmapsto A+\eps_i(2n+k),\qquad
 T\longmapsto T+\eps_iR_{r,n}
\]
preserve the exponent pairs, hence \(h\), \(v\), and the estimate \eqref{eq:semigroup-cone-bound}.  They commute with extraction of the \(A\)-constant term, with the induced substitution \(Q\mapsto Qe^{-\eps_iR_{r,n}}\) on the target:
\[
 \operatorname{CT}_A(\tau_{i,n}f)
 =\tau_{i,n}^{\,0}(\operatorname{CT}_A f).
\]
Here \(\tau_{i,n}^{\,0}\) denotes this \(Q\)-substitution together with the corresponding change of equivariant parameters.  This is commutation, not invariance of the resulting \(Q\)-series.  The common monoid includes the generators for all three parameter pairs before the additive translations are applied.  The factor
\[
 \left(-\sqrt{tq}\,\frac uQ;t,q\right)_\infty
\]
lies in \(1+\mathfrak m_A\).  Under Conjecture~\ref{conj:P2-asymptotic-resummation}, the specified expansions of \(\mathcal R\) and every chart translate lie in \(1+\mathfrak m_A\).  Dividing by the corresponding incoming product shows that the induced expansions of \(\Phi(Y_{\mathrm{edge}})\Phi(Y_{\mathrm{root}})^{-1}\) and its chart translates do so as well.
\end{lemma}

\begin{equation}
 \operatorname{CT}_A\bigl(
 \mathscr A_{\mathrm c}(Q,u;t,q)\mathcal R(A,T;\eps_1,\eps_2)\bigr)
 =\mathscr A_{\mathrm c}(Q,0;t,q)=\mathcal Z(Q;t,q).
\label{eq:u-zero-order-zero}
\end{equation}

\begin{lemma}
\label{lem:two-chart-factor}
For the chart substitutions of Lemma~\ref{lem:NY-chart-substitutions}, let
\begin{equation}
 A_i=A+\eps_i(2n+k),\qquad
 M_i=M+\eps_i\eta_{k,d_{\mathrm{NY}}},\qquad
 \eta_{k,d_{\mathrm{NY}}}
 =\frac{3-2d_{\mathrm{NY}}-k}{2}.
\label{eq:centered-additive-chart-data}
\end{equation}
The quotient of the exponentials of \(G_{\alpha,\beta}(A,M)+\widetilde P_M(M;\alpha,\beta)\) at the two chart substitutions by the unshifted exponential is
\begin{equation}
 (-1)^n e^{-\pi\mathrm{i}k/2}
 \exp\left[
 \left(\frac1{12}-\frac{\eta_{k,d_{\mathrm{NY}}}^2}{3}\right)M
 -s\left(\frac{\eta_{k,d_{\mathrm{NY}}}^3}{9}
          +\frac1{24}\right)\right].
\label{eq:two-chart-factor}
\end{equation}
\end{lemma}

\begin{equation}
 \exp\left(\frac{M}{12}-\frac{s}{24}\right).
\label{eq:zero-common-M-factor}
\end{equation}

\begin{lemma}
\label{lem:bilateral-u-zero-limit}
Assume Conjectures~\ref{conj:P2-vertex-comparison} and \ref{conj:P2-asymptotic-resummation}, and fix the angular region \(\mathfrak S\) above.  For
\[
 (k,d_{\mathrm{NY}})=(0,1),\qquad (0,2),\qquad (1,1),
\]
remove the factors of Lemma~\ref{lem:two-chart-factor} from the bilateral sum over \(n\in\Z\); in the last case also remove \(Q^{1/3}\), the factor \eqref{eq:zero-common-M-factor}, and the constants described in the proof of that lemma.  The resulting sums lie in \(\mathscr H_{\mathfrak S}\), and \(\operatorname{CT}_A\) may be applied term by term.
\end{lemma}

\begin{proof}
Substituting \(R=3n+\tfrac12,3n-\tfrac12,3n+\tfrac32\) in \eqref{eq:P2-degree-zero-polynomial} gives the initial \(Q\)-exponents
\[
 m_1(n)
 \quad\text{for }(k,d_{\mathrm{NY}})=(0,1),\qquad
 m_{-1}(n)
 \quad\text{for }(k,d_{\mathrm{NY}})=(0,2),
\]
and
\[
 \frac13+m_3(n)
 \quad\text{for }(k,d_{\mathrm{NY}})=(1,1).
\]
In the last case, first remove \(Q^{1/3}\).  By Lemma~\ref{lem:completed-semigroup-support}, every translated \(\mathcal R\)-factor and every translated \(\mathscr A_{\mathrm c}\)-factor lies in \(\mathscr H_{\mathfrak S}\), and all three functions \(m_r\) have finite sublevel sets on \(\Z\).  Let \(\mathbf e^\gamma\) be a monomial of the product of the two translated \(\mathcal R\)-factors, with \(v(\gamma)=\delta\), and write the two monomials selected from the translated \(\mathscr A_{\mathrm c}\)-factors as \(Q^{a_1}u^{b_1}\) and \(Q^{a_2}u^{b_2}\), \(a_i,b_i\geq0\); their chart translations change only the coefficients.  A term with total \(Q\)-exponent \(d_{\mathrm{tot}}\) satisfies
\[
 m_r(n)+a_1+a_2+\delta=d_{\mathrm{tot}},
\]
and since \(\delta\geq-Ch(\gamma)\),
\[
 m_r(n)\leq d_{\mathrm{tot}}+Ch(\gamma).
\]
To check the support condition of Definition~\ref{def:completed-semigroup-algebra}, fix bounds \(h\leq H\) and \(v\leq V_0\).  The \(h\)-value of the term is
\[
 h(\gamma)+(b_1+b_2)h(1,0)
\]
and its \(v\)-value is
\[
 v(\gamma)+m_r(n)+a_1+a_2.
\]
Since \(h(1,0)>0\), the first bound leaves finitely many \(b_1+b_2\), hence finitely many \((b_1,b_2)\).  It also gives \(h(\gamma)\leq H\), while the second gives \(v(\gamma)\leq V_0\) since \(m_r(n),a_1,a_2\geq0\), and the common finitely generated monoid has only finitely many such \(\gamma\), independently of \(n\): the \(h\)-bound bounds all positive-\(h\) generator multiplicities, and the \(v\)-bound then bounds the multiplicity of \((0,1)\).  For each \(\gamma\), \eqref{eq:semigroup-cone-bound} gives
\[
 m_r(n)\leq V_0-v(\gamma)\leq V_0+CH.
\]
Only finitely many \(n\) occur, and once \(\gamma,n,b_1,b_2\) are fixed, the inequality \(a_1+a_2\leq V_0-v(\gamma)-m_r(n)\) leaves finitely many \((a_1,a_2)\).  The bilateral sum therefore has locally finite support and lies in \(\mathscr H_{\mathfrak S}\).  This is a coefficientwise sum of a locally finite family; it need not be the limit of its partial sums in the \(h\)-filtration, since the \(Q\)-series have \(h=0\).

By Lemma~\ref{lem:completed-semigroup-support}, the chart translations preserve \(h,v\) and commute with \(\operatorname{CT}_A\) with the induced \(Q\)-substitution.  The coefficientwise finiteness just proved allows \(\operatorname{CT}_A\) to be applied term by term; its \(h=0\) part selects the constant term \(1\) of each \(\mathcal R\)-factor and the constant term in \(u\) of each \(\mathscr A_{\mathrm c}\)-factor, which is the shifted series \(\mathcal Z\) by \eqref{eq:u-zero-order-zero}.
\end{proof}

\end{document}
